\documentclass[14pt,a4paper]{extarticle}

% ---------- Кодировка и язык ----------
\usepackage[utf8]{inputenc}
\usepackage[T2A]{fontenc}
\usepackage[russian]{babel}

% ---------- Поля страницы ----------
\usepackage{geometry}
\geometry{
    left=30mm,
    right=10mm,
    top=20mm,
    bottom=20mm
}

% ---------- Шрифт и интервалы ----------
\usepackage{tempora} % аналог Times New Roman
\usepackage{setspace}
\onehalfspacing

% ---------- Математика ----------
\usepackage{amsmath, amssymb, amsthm}

% ---------- Рисунки и таблицы ----------
\usepackage{graphicx}
\usepackage{float}
\usepackage{caption}
\usepackage{subcaption}
\usepackage{array}
\usepackage{booktabs}

% ---------- Абзацы ----------
\usepackage{indentfirst}
\setlength{\parindent}{1.25cm}
\setlength{\parskip}{0pt}

% ---------- Ссылки ----------
\usepackage{hyperref}
\hypersetup{
    colorlinks=false,
    pdfborder={0 0 0}
}

% ---------- Нумерация формул, таблиц, рисунков ----------
\numberwithin{equation}{section}
\numberwithin{figure}{section}
\numberwithin{table}{section}

% ---------- Оформление заголовков ----------
\usepackage{titlesec}

\titleformat{\section}
  {\normalfont\bfseries}
  {\thesection}
  {1em}
  {}

\titleformat{\subsection}
  {\normalfont\bfseries}
  {\thesubsection}
  {1em}
  {}

\titleformat{\subsubsection}
  {\normalfont\bfseries}
  {\thesubsubsection}
  {1em}
  {}

% ---------- Подписи ----------
\captionsetup[figure]{
    labelsep=endash,
    justification=centering
}

\captionsetup[table]{
    labelsep=endash,
    justification=raggedright,
    singlelinecheck=false
}

\titlespacing*{\section}
{1.25cm}{*3}{*2}

\titlespacing*{\subsection}
{1.25cm}{*2}{*1.5}

\titlespacing*{\subsubsection}
{1.25cm}{*2}{*1}


% ---------- Документ ----------
\begin{document}

% ============================================================
% СОДЕРЖАНИЕ
% ============================================================

\tableofcontents
\newpage

% ============================================================
% ВВЕДЕНИЕ
% ============================================================

\section*{Введение}
\addcontentsline{toc}{section}{Введение}

В настоящее время активно развиваются сенсоры и устройства, которые способны детектировать малые механические воздействия окружающей среды. Существующие в природе структуры и механизмы восприятия, такие как: реснички, волоски насекомых, латеральной линии рыб демонстрируют высокую чувствительностью к микроскопическим силам, вибрациям, потокам жидкостей или газов, благодаря чему они находят широкое применение в разработке детекторов мониторинга окружающей среды, робототехнике, медицинской диагностики \cite{ribeiro2017,man2023}.

Одним из перспективных направлений является разработка магнитных сенсоров, в основе которых композиты полимерные структуры и магнитных частиц. Принцип работы подобных устройств заключается в фиксации магнитным датчиком, например датчиком Холла или магнитным сенсором, изменения магнитного поля, вызванного деформацией композитного элемента структуры под действием внешних механических сил. Таким образом, механическая величина может быть восстановлена по изменению магнитного сигнала\cite{ribeiro2017,man2023}.

В настоящей работе был рассмотрен датчик, описанный в статье A tactile and airflow motion sensor based on flexible double-layer magnetic cilia \cite{man2023}. В данной работе предложена двухслойная магнитная ресничка, нижний слой которой выполнен из чистого полидиметилсилоксана (PDMS), верхний --- из композита PDMS с магнитными частицами NdFeB. Подобная структура позволяет совместить два полезных свойства: высокую гибкость, благодаря нижнему слою и достаточную намагниченность композитного. Авторами публикации показано, что описанный сенсор способен измерять механические силы в диапазоне $0$--$60~\mu N$ с разрешением $2.1~\mu N$, а также регистрировать воздушные потоки с чувствительностью $1.43~\mu T/(\mathrm{m/s})$ \cite{man2023}.

Для описания работы сенсора необходимо рассматривать две взаимосвязанные задачи: механическую задачу деформирования мягкой двухслойной реснички и задачу расчёта магнитного поля, создаваемого магнитными включениями композита.

Актуальность данной работы связана с необходимостью численного
моделирования таких сенсорных систем ещё на этапе проектирования.
Экспериментальное изготовление и испытание магнитных микроструктур
требует значительных технологических затрат, тогда как математическое и
компьютерное моделирование позволяет заранее оценивать влияние
геометрии, механических свойств материала, распределения магнитных
частиц и величины внешнего воздействия на чувствительность датчика.
Для задач мягких магнитных композитов и магнитоуправляемых структур
широко применяются методы нелинейной механики и численного
моделирования, в частности метод конечных элементов \cite{zhao2019,kim2022}.


\newpage

% ============================================================
% ГЛАВА 1
% ============================================================

\section{Обзор литературы}

\subsection{Метод конечных элементов}
Метод конечных элементов на данный момент является одним из основных численных методов для решения краевых задач математической физики. Метод основывается на представлении сложной непрерывной области как конечного количества простых субобластей (конечных элементов).  Внутри каждого из элементов искомое поле, то есть функция, заданная в каждой точке выбранной области, приближенно представляется через набор простых, чаще всего полиномиальных, базисных функций. По итогу сложная исходная задача решения дифференциального уравнения в частных производных сводится к системе алгебраических уравнений конечного числа неизвестных коэффициентов, представляющих собой значения искомого поля в узлах сетки.\cite{Zienkiewicz2005}

Идеи, впоследствии примененные в реализации метода, впервые упоминаются в работe Р. Куранта 1943 года, в которой автор использовал вариационный подход и кусочно-линейную аппроксимацию для решения задачи кручения\cite{courant1943}. В современном виде, как метод решения инженерных прикладных задач, МКЭ обрел популярность в 1950 году. Ярким примером является работа \cite{turner1956}, посвященная расчету жесткости и деформации сложных конструкций. Впоследствии метод неоднократно модернизировался, потому сейчас применяется в широком диапазоне прикладных задач, став одним из базовых инструментов вычислительной механики, механики деформируемого твердого тела, гидродинамики, теплопроводности, электромагнетизма и других мультифизических задач \cite{Zienkiewicz2005, brenner2008}.

Математически метод конечных элементов заключается в переходе от сильной дифференциальной постановки задачи к слабой (вариационной) постановке.

Пусть в области $\Omega \subset \mathbb{R}^n$ требуется найти функцию $u$,
удовлетворяющую уравнению с заданными граничными условиями:

\begin{equation}
    \label{eq:mean_number}
    \mathcal{L}(u)=f \quad \text{в } \Omega,
\end{equation}

После умножения на тестовую функцию
$v$ и интегрирования по области получают интегральное соотношение. Для
многих задач после интегрирования по частям оно принимает вид:

\begin{equation}
    \label{eq:mean_number}
    a(u,v)=l(v),
\end{equation}

где $a(u,v)$ --- билинейная форма; \\
$l(v)$ --- линейный функционал.

Искомое решение затем аппроксимируется функцией из конечномерного
пространства:

\begin{equation}
    \label{eq:mean_number}
    u_h(x)=\sum_{i=1}^{N} U_i \varphi_i(x),
\end{equation}

где $\varphi_i(x)$ --- базисные функции; \\
$U_i$ --- неизвестные коэффициенты.

Подстановка этого представления в слабую форму приводит к
системе алгебраических уравнений

\begin{equation}
    \label{eq:mean_number}
    \mathbf{K}\mathbf{U}=\mathbf{F}.
\end{equation} \\

В задачах нелинейной механики вместо линейной системы возникает
нелинейная система \\

\begin{equation}
    \label{eq:mean_number}
    \mathbf{R}(\mathbf{U})=\mathbf{0},
\end{equation} \\

Которая обычно решается итерационными методами, например методом Ньютона: \\

\begin{equation}
    \label{eq:mean_number}
    \mathbf{K}(\mathbf{U}_k)\Delta \mathbf{U}_k
    =
    -\mathbf{R}(\mathbf{U}_k),
\end{equation} \\

\begin{equation}
    \label{eq:mean_number}
    \mathbf{U}_{k+1}
    =
    \mathbf{U}_k+\Delta\mathbf{U}_k.
\end{equation} \\


\subsection{Общий принцип метода конечных элементов}

Метод конечных элементов (МКЭ) является численным методом решения краевых задач математической физики, основанным на аппроксимации искомой функции в конечномерном пространстве. \cite{Zienkiewicz2005}

Пусть требуется решить краевую задачу в области
\[
\Omega \subset \mathbb{R}^n
\]
в виде дифференциального уравнения:
\[
\mathcal{L}(u) = f \quad \text{в } \Omega,
\]
с граничными условиями:
\[
u = g \quad \text{на } \Gamma_D, \qquad
\mathcal{B}(u) = h \quad \text{на } \Gamma_N,
\]
где:
\begin{itemize}
    \item $u$ — искомая функция,
    \item $\mathcal{L}$ — дифференциальный оператор,
    \item $f$ — заданная правая часть,
    \item $\Gamma_D$ — часть границы с условиями Дирихле,
    \item $\Gamma_N$ — часть границы с естественными условиями.
\end{itemize}

\subsection{Переход к слабой (вариационной) форме}

Основная идея МКЭ заключается в переходе от дифференциальной (сильной) формы к интегральной (слабой) форме.

Вводится пространство тестовых функций $v$, и исходное уравнение умножается на $v$ и интегрируется по области:
\[
\int_{\Omega} \mathcal{L}(u)\, v \, d\Omega = \int_{\Omega} f v \, d\Omega.
\]

После интегрирования по частям (снижения порядка производных) получаем слабую форму:
\[
a(u, v) = l(v),
\]
где:
\begin{itemize}
    \item $a(u,v)$ — билинейная форма,
    \item $l(v)$ — линейный функционал.
\end{itemize}

Задача формулируется следующим образом:

\textit{Найти $u \in V$, такое что}
\[
a(u, v) = l(v), \quad \forall v \in V_0,
\]
где $V$ — функциональное пространство допустимых решений.

\subsection{Разбиение области (дискретизация)}

Область $\Omega$ разбивается на конечное число непересекающихся подобластей (конечных элементов):
\[
\Omega_h = \bigcup_{K \in \mathcal{T}_h} K,
\]
где $\mathcal{T}_h$ — сетка.

В трёхмерном случае используются тетраэдры:
\[
K \subset \mathbb{R}^3.
\]

Характерный размер элемента обозначается через $h$.

\subsection{Аппроксимация решения}

Искомая функция $u$ аппроксимируется в конечномерном пространстве:
\[
u_h(x) = \sum_{i=1}^{N} U_i \varphi_i(x),
\]
где:
\begin{itemize}
    \item $\varphi_i$ — базисные функции (локальные полиномы),
    \item $U_i$ — неизвестные коэффициенты (степени свободы),
    \item $N$ — число узлов сетки.
\end{itemize}

Для линейных элементов:
\[
\varphi_i \in \mathbb{P}_1,
\]
то есть функции линейны внутри каждого элемента.

\subsection{Формирование системы уравнений}

Подставляя аппроксимацию $u_h$ в слабую форму, получаем систему алгебраических уравнений:
\[
\mathbf{K}\mathbf{U} = \mathbf{F},
\]
где:
\begin{itemize}
    \item $\mathbf{K}$ — матрица жёсткости,
    \item $\mathbf{U}$ — вектор неизвестных,
    \item $\mathbf{F}$ — вектор нагрузок.
\end{itemize}

Элементы матрицы:
\[
K_{ij} = a(\varphi_j, \varphi_i),
\]
\[
F_i = l(\varphi_i).
\]

\subsection{Нелинейные задачи}

В задачах нелинейной механики оператор $\mathcal{L}$ зависит от решения $u$:
\[
\mathcal{L}(u) = 0.
\]

После дискретизации получается нелинейная система:
\[
\mathbf{R}(\mathbf{U}) = 0.
\]

Для её решения применяется метод Ньютона:
\[
\mathbf{K}(\mathbf{U}_k)\Delta \mathbf{U}_k = -\mathbf{R}(\mathbf{U}_k),
\]
\[
\mathbf{U}_{k+1} = \mathbf{U}_k + \Delta \mathbf{U}_k.
\]

Здесь:
\begin{itemize}
    \item $\mathbf{R}$ — вектор невязки,
    \item $\mathbf{K}$ — касательная матрица (якобиан).
\end{itemize}

\subsection{Алгоритм метода конечных элементов}

Общий алгоритм МКЭ включает следующие этапы:

\begin{enumerate}
    \item Построение геометрии области $\Omega$;
    \item Разбиение области на конечные элементы (сетка);
    \item Выбор базисных функций;
    \item Формирование слабой формы задачи;
    \item Сборка глобальной системы уравнений;
    \item Задание граничных условий;
    \item Решение линейной или нелинейной системы;
    \item Постобработка результатов.
\end{enumerate}

\subsection{Сходимость метода}

При уменьшении размера элемента $h \to 0$ решение МКЭ стремится к точному решению:
\[
u_h \to u.
\]

Погрешность аппроксимации оценивается как:
\[
\|u - u_h\| \le C h^p,
\]
где:
\begin{itemize}
    \item $p$ — порядок аппроксимации,
    \item $C$ — константа.
\end{itemize}


\subsection{Магнитостатическое моделирование и конечно-элементные методы в задачах магнитных сенсоров}

Моделирование магнитных сенсоров, выполненых на основе мягких магнитных композитов, 
учитывает две физические составляющие: механическую деформацию гибкого элемента, а также изменение 
магнитного поля, вызванное смещением и поворотом намагниченной области. Например, в работе \cite{man2023},
в которой предложен тактильный и аэродинамический датчик.
Чувствителеный элемент сенсора представляет собой гибридный материал, нижний сегмент которого 
выполнен из мякого полидиметилсилоксана (PDMS), а верхний из композита PDMS и магинтных частиц NdFeB,
яаляющихся источником магнитного поля. Авторы указывают, что при изгибе ресничи, направление магниного поля
внутри датчка Холла изменяется с вертикального на наклонное, что позволяет связать механическое воздейсвтеи
с эдектрическм выходным сигналом сенсора Холла. 

С точки зрения физики магнитостатическая задача является частным случаем уравнений Максвелла, в которых 
можно пренебречь временными производными. В классической постановке магнитостатическое поле описывается
следующим уравнениями

\[
    \nabla \times \mathbf H = \mathbf J,
\]
\[
    \nabla \cdot \mathbf B = 0,
\]
где $\mathbf H$ --- напряжённость магнитного поля;

$\mathbf B$ --- магнитная индукция; 

$\mathbf J$ --- плотность свободного электрического тока. Связь между магнитной индукцией, напряжённостью.

\[
    \mathbf B = \mu_0(\mathbf H + \mathbf M),
\]

где $\mu_0$ --- магнитная постоянная;

$\mathbf M$ --- вектор намагниченности.

Такая постановка магнитостатической задачи подробно рассматривается в работах по вычислительной 
электромагнетике и конечно-элементнму моделированию \cite{Bossavit1998,Jin2014,Kuczmann2008,Bermudez2008}.
Наприемр, авторы работы \cite{Bermudez2008} формулируют  трехмерную магнитостатическую задачу в 
ограниченной области, содержащей токи и магнитные материалы.

В задачх без свободных токов:

\[
    \mathbf J = 0.
\]

Тогда из уравнения

\[
    \nabla \times \mathbf H = 0
\]

следует, что в односвязной области напряжённость магнитного поля может быть представлена через 
скалярный магнитный потенциал $\varphi$:

\[
    \mathbf H = -\nabla \varphi.
\]

Подставляя это выражение в материальное соотношение, получаем

\[
    \mathbf B = \mu_0(-\nabla \varphi + \mathbf M).
\]

Далее используется условие отсутствия магнитных зарядов:

\[
    \nabla \cdot \mathbf B = 0.
\]

В результате получается уравнение для скалярного магнитного потенциала:

\[
    \nabla \cdot (-\nabla \varphi + \mathbf M) = 0.
\]

Для решения методом конечных элементов данное уравнение переводится в слабую форму. То есть, потенциал
ищется в функциональном пространстве, уравнение умножается на тестовую функцию и интегрируется по
расчетной области. 

\[
    \int_{\Omega} \nabla \varphi \cdot \nabla v \, d\Omega
    =
    \int_{\Omega_m} \mathbf M \cdot \nabla v \, d\Omega,
\]
где $\Omega$ --- расчётная область, содержащая воздух и магнитный материал;

$\Omega_m$ --- область намагниченного композитного слоя.


Искомый потенциал аппроксимируется в конечномерном пространстве:

\[
    \varphi_h(\mathbf x)
    =
    \sum_{i=1}^{N} \Phi_i N_i(\mathbf x),
\]
где $N_i(\mathbf x)$ --- базисные функции конечных элементов;

$\Phi_i$ --- неизвестные узловые значения потенциала.
 
После подстановки этого представления в слабую форму задача сводится к системе линейных
алгебраических уравнений, которые представляют общей основой метода конечных элементов, описанной в
\cite{Zienkiewicz2005}:

\[
    K \Phi = F.
\]

Элементы матрицы и правой части имеют вид:

\[
    K_{ij}
    =
    \int_{\Omega} \nabla N_j \cdot \nabla N_i \, d\Omega,
\]
\[
    F_i
    =
    \int_{\Omega_m} \mathbf M \cdot \nabla N_i \, d\Omega.
\]

После решения системы магнитная индукция в воздушной области восстанавливается по формуле

\[
    \mathbf B_{\text{air}}
    =
    -\mu_0 \nabla \varphi_h.
\]

Скалярно-потенциальная постановка является не единственным способом конечно-элементного моделирования магнитостатики.
В задачах с токами часто используется векторный магнитный потенциал $\mathbf A$, определяемый соотношением:

\[
    \mathbf B = \nabla \times \mathbf A.
\]

Такое представление автоматически удовлетворяет условию

\[
    \nabla \cdot \mathbf B = 0.
\]

Однако векторно-потенциальная постановка требует использования векторных конечных элементов и дополнительной калибровки потенциала,
поскольку $\mathbf A$ определяется неоднозначно. В задачах без свободных токов, когда основной источник поля задаётся намагниченностью $\mathbf M$,
более экономичной является постановка через скалярный потенциал. Различные потенциальные формулировки магнитостатических задач и
их конечно-элементная реализация подробно рассматриваются в работах \cite{Bossavit1998, Jin2014, Kuczmann2008, Bermudez2008}.

Существенной особенностью магнитостатического моделирования является необходимость задания внешней границы расчётной области.
В реальной физической задаче магнитное поле существует во всём пространстве, то есть во внешней бесконечной области.
При численном решении расчётная область неизбежно ограничивается конечным объёмом. Поэтому внешняя граница является искусственной,
а результат может зависеть от её положения и от выбранного граничного условия. Если на внешней границе задаётся условие $\varphi = 0$,
то это соответствует грубому переносу условия затухания потенциала на конечное расстояние. Такая аппроксимация может искажать поле,
если граница расположена слишком близко к магнитному источнику. Альтернативой является использование естественного граничного условия,
однако оно также не является точным описанием бесконечной области. Поэтому в магнитостатических FEM-расчётах необходима проверка 
чувствительности решения к размерам воздушной области.

Отдельный класс работ посвящён мягким магнитным материалам, в которых магнитные частицы встроены в деформируемую полимерную матрицу. В работе 
\cite{Zhao2019} разработана теория hard-magnetic soft materials, описывающая связь между конечными деформациями мягкого тела и магнитными воздействиями.
Авторы построили конститутивную модель, реализовали её в конечно-элементной среде и показали, что такая модель позволяет описывать большие деформации 
магнитомягких структур под действием магнитного поля.

Обзор \cite{Kim2022} систематизирует современные подходы к магнитным мягким материалам и роботам. В нём рассмотрены способы изготовления магнитных композитов, 
методы программирования намагниченности, механизмы магнитного управления и области применения таких материалов. Работа показывает общий контекст использования 
мягких магнитных композитов: изменение формы магнитной структуры приводит к изменению создаваемого ею поля.

Магнитостатическое моделирование также широко используется при проектировании мягких тактильных сенсоров на основе датчиков Холла. 
В работе \cite{Wang2016} предложена методология проектирования мягкого трёхосевого тактильного сенсора, в котором внешняя сила вызывает 
смещение магнита относительно Hall-сенсора. Авторы использовали численное моделирование для анализа связи между приложенной силой, 
деформацией мягкой структуры и изменением компонент магнитного поля. В работе \cite{DeBoer2017} конечно-элементное моделирование применялось 
для оптимизации магнитного мягкого тактильного сенсора. Авторы исследовали влияние геометрических параметров и механических свойств эластомерного 
слоя на чувствительность устройства. В работе \cite{Rosle2019} похожий подход использовался для оптимизации геометрии сенсора, что позволяло увеличить 
информативность магнитного сигнала при контакте с объектом.


\newpage
% ============================================================
% ГЛАВА 2
% ============================================================
\section{Математическая постановка задачи механического деформирования реснички}

\subsection{Описание расчетной области}

В данной работе рассматривается трехмерная модель двухслойной магнитной реснички,
закрепленной на цилиндрической подложке. Такая конструкция соответствует
чувствительному элементу датчика, описанного в работе \cite{man2023}: нижний слой
реснички выполнен из чистого PDMS, а верхний слой представляет собой композит
PDMS с магнитными частицами NdFeB.

Полная расчетная область имеет вид
\[
    \Omega =
    \Omega_s \cup \Omega_1 \cup \Omega_2,
\]
где
\(\Omega_s\) --- область подложки,
\(\Omega_1\) --- нижний слой реснички из PDMS,
\(\Omega_2\) --- верхний магнитный композитный слой.

Ось реснички направлена вдоль оси \(z\). Радиус реснички равен
\[
    R = \frac{D}{2}.
\]
В численной модели используются следующие геометрические параметры:
\[
    D = 120 \cdot 10^{-6}\,\text{м},
\]
\[
    L_1 = 2 \cdot 10^{-3}\,\text{м},
    \qquad
    L_2 = 2 \cdot 10^{-3}\,\text{м}.
\]
Полная длина реснички равна
\[
    L = L_1 + L_2.
\]

Подложка моделируется цилиндром радиуса \(R_s\) и толщины \(H_s\):
\[
    R_s = 0.60 \cdot 10^{-3}\,\text{м},
    \qquad
    H_s = 0.50 \cdot 10^{-3}\,\text{м}.
\]

Таким образом, области задаются следующим образом:
\[
    \Omega_s =
    \left\{
    (x,y,z): x^2+y^2 \le R_s^2,\;
    -H_s \le z \le 0
    \right\},
\]
\[
    \Omega_1 =
    \left\{
    (x,y,z): x^2+y^2 \le R^2,\;
    0 \le z \le L_1
    \right\},
\]
\[
    \Omega_2 =
    \left\{
    (x,y,z): x^2+y^2 \le R^2,\;
    L_1 \le z \le L
    \right\}.
\]

\subsection{Неизвестная функция}

Искомой величиной является вектор перемещений
\[
    \mathbf{u}(\mathbf{X}) =
    \begin{pmatrix}
        u_x(\mathbf{X})\\
        u_y(\mathbf{X})\\
        u_z(\mathbf{X})
    \end{pmatrix},
    \qquad
    \mathbf{X}\in\Omega.
\]
Здесь \(\mathbf{X}\) обозначает координату точки в исходной, недеформированной
конфигурации.

После деформации точка \(\mathbf{X}\) переходит в положение
\[
    \mathbf{x} = \boldsymbol{\varphi}(\mathbf{X})
    =
    \mathbf{X} + \mathbf{u}(\mathbf{X}).
\]

\subsection{Кинематика конечных деформаций}

Поскольку ресничка является мягкой и может испытывать большие перемещения,
в работе используется нелинейная постановка механики деформируемого тела.

Градиент деформации определяется формулой
\[
    \mathbf{F}
    =
    \nabla_{\mathbf{X}}\boldsymbol{\varphi}
    =
    \mathbf{I} + \nabla_{\mathbf{X}}\mathbf{u},
\]
где \(\mathbf{I}\) --- единичный тензор.

Правый тензор Коши--Грина равен
\[
    \mathbf{C} = \mathbf{F}^{T}\mathbf{F}.
\]

Первый инвариант тензора \(\mathbf{C}\) имеет вид
\[
    I_C = \operatorname{tr}\mathbf{C}.
\]

Изменение объема описывается якобианом деформации
\[
    J = \det \mathbf{F}.
\]
При \(J=1\) локальный объем материала сохраняется. Значение \(J\neq 1\)
соответствует локальному изменению объема.

\subsection{Материальная модель}

Материалы считаются гиперупругими и описываются сжимаемой моделью
Neo-Hookean. Плотность энергии деформации задается выражением
\[
    \psi(\mathbf{F})
    =
    \frac{\mu}{2}
    \left(
        \operatorname{tr}(\mathbf{C}) - 3
    \right)
    -
    \mu \ln J
    +
    \frac{\lambda}{2}(\ln J)^2.
\]

Параметры Ламе выражаются через модуль Юнга \(E\) и коэффициент Пуассона
\(\nu\):
\[
    \mu = \frac{E}{2(1+\nu)},
\]
\[
    \lambda =
    \frac{E\nu}{(1+\nu)(1-2\nu)}.
\]

В численной модели используются кусочно-постоянные материальные параметры:
\[
    E(\mathbf{X}) =
    \begin{cases}
        E_{\mathrm{PDMS}}, & \mathbf{X}\in \Omega_s \cup \Omega_1,\\
        E_{\mathrm{mag}}, & \mathbf{X}\in \Omega_2,
    \end{cases}
\]
\[
    \nu(\mathbf{X}) =
    \begin{cases}
        \nu_{\mathrm{PDMS}}, & \mathbf{X}\in \Omega_s \cup \Omega_1,\\
        \nu_{\mathrm{mag}}, & \mathbf{X}\in \Omega_2.
    \end{cases}
\]

Основные значения параметров:
\[
    E_{\mathrm{PDMS}} = 1.5\cdot 10^{6}\,\text{Па},
\]
\[
    E_{\mathrm{mag}} = 16.6\cdot 10^{6}\,\text{Па}.
\]

В базовой серии расчетов принимается
\[
    \nu_{\mathrm{PDMS}}=\nu_{\mathrm{mag}}=0.49.
\]
Дополнительно проводится диагностический расчет при
\[
    \nu_{\mathrm{PDMS}}=\nu_{\mathrm{mag}}=0.45
\]
для проверки чувствительности модели к почти несжимаемости материала и
возможному эффекту объемного запирания.

\subsection{Тензор напряжений}

Первый тензор Пиолы--Кирхгофа определяется как производная плотности
энергии по градиенту деформации:
\[
    \mathbf{P}
    =
    \frac{\partial \psi}{\partial \mathbf{F}}.
\]

Для выбранной модели Neo-Hookean он имеет вид
\[
    \mathbf{P}
    =
    \mu
    \left(
        \mathbf{F} - \mathbf{F}^{-T}
    \right)
    +
    \lambda \ln J\, \mathbf{F}^{-T}.
\]

\subsection{Сильная форма задачи равновесия}

Внутри тела объемные силы не учитываются. Поэтому уравнение равновесия
в исходной конфигурации записывается в виде
\[
    \operatorname{Div}\mathbf{P} = \mathbf{0}
    \qquad
    \text{в } \Omega.
\]

Здесь \(\operatorname{Div}\) обозначает дивергенцию по материальным
координатам.

\subsection{Граничные условия}

Нижняя поверхность подложки жестко закреплена:
\[
    \Gamma_{\mathrm{bottom}}
    =
    \left\{
    (x,y,z)\in\partial\Omega:
    z=-H_s
    \right\}.
\]
На этой границе задается условие Дирихле
\[
    \mathbf{u} = \mathbf{0}
    \qquad
    \text{на } \Gamma_{\mathrm{bottom}}.
\]

В актуальной версии модели управляющее смещение задается не на боковой
контактной области, а на всей верхней грани реснички:
\[
    \Gamma_{\mathrm{top}}
    =
    \left\{
    (x,y,z)\in\partial\Omega:
    z=L
    \right\}.
\]

На этой границе задается горизонтальное смещение по оси \(x\):
\[
    u_x = \delta_x
    \qquad
    \text{на } \Gamma_{\mathrm{top}}.
\]

Компоненты \(u_y\) и \(u_z\) на верхней грани явно не фиксируются:
\[
    u_y,\;u_z
    \quad
    \text{остаются свободными на } \Gamma_{\mathrm{top}}.
\]

На остальных частях границы внешняя поверхностная нагрузка отсутствует,
поэтому выполняется естественное граничное условие
\[
    \mathbf{P}\mathbf{N}=\mathbf{0}
    \qquad
    \text{на } \Gamma_{\mathrm{free}},
\]
где \(\mathbf{N}\) --- внешняя нормаль к границе в исходной конфигурации.

\subsection{Вариационная постановка}

Полная энергия деформации тела задается функционалом
\[
    \Pi(\mathbf{u})
    =
    \int_{\Omega}
    \psi(\mathbf{F}(\mathbf{u}))\,dV.
\]

Равновесное состояние определяется условием стационарности полной энергии:
\[
    \delta \Pi(\mathbf{u};\mathbf{v}) = 0
\]
для всех допустимых виртуальных перемещений \(\mathbf{v}\), удовлетворяющих
однородным условиям на границах с заданными перемещениями.

Слабая форма задачи имеет вид:
\[
    \boxed{
    \int_{\Omega}
    \mathbf{P}(\mathbf{u}):\nabla \mathbf{v}\,dV = 0
    }
\]
для всех допустимых тестовых функций \(\mathbf{v}\).

Здесь двоеточие означает двойное скалярное произведение тензоров:
\[
    \mathbf{A}:\mathbf{B}
    =
    \sum_{i,j} A_{ij}B_{ij}.
\]

\subsection{Конечно-элементная аппроксимация}

Расчетная область разбивается на тетраэдральную сетку
\[
    \Omega_h = \bigcup_{K\in\mathcal{T}_h} K,
\]
где \(K\) --- тетраэдральный конечный элемент, а \(h\) --- характерный
размер элемента.

В базовых расчетах используется сетка с
\[
    h = 50\cdot 10^{-6}\,\text{м}.
\]

Вектор перемещений аппроксимируется в конечно-элементном пространстве
лагранжевых элементов степени \(p\):
\[
    \mathbf{u}_h(\mathbf{X})
    =
    \sum_{i=1}^{N}
    \mathbf{U}_i \varphi_i(\mathbf{X}),
\]
где
\(\varphi_i\) --- базисные функции,
\(\mathbf{U}_i\) --- неизвестные узловые векторы перемещений.

В расчетах рассматриваются элементы первого и второго порядка:
\[
    p=1
    \qquad
    \text{и}
    \qquad
    p=2.
\]

Материальные параметры \(E\) и \(\nu\) задаются в пространстве \(DG0\), то есть
являются постоянными внутри каждого тетраэдрального элемента и могут
скачком изменяться на границе слоев.

\subsection{Дискретная нелинейная система}

После подстановки конечно-элементного приближения \(\mathbf{u}_h\) в слабую
форму получается нелинейная система алгебраических уравнений:
\[
    \mathbf{R}(\mathbf{U})=\mathbf{0}.
\]

Здесь \(\mathbf{U}\) --- глобальный вектор неизвестных узловых перемещений:
\[
    \mathbf{U}
    =
    \left(
    U_{1x},U_{1y},U_{1z},
    \ldots,
    U_{Nx},U_{Ny},U_{Nz}
    \right)^T.
\]

Для решения нелинейной системы используется метод Ньютона. На \(k\)-й
итерации решается линейная система
\[
    \mathbf{K}(\mathbf{U}_k)\Delta \mathbf{U}_k
    =
    -\mathbf{R}(\mathbf{U}_k),
\]
где
\[
    \mathbf{K}(\mathbf{U}_k)
    =
    \frac{\partial \mathbf{R}}{\partial \mathbf{U}}(\mathbf{U}_k)
\]
--- касательная матрица жесткости.

После решения линейной системы выполняется обновление:
\[
    \mathbf{U}_{k+1}
    =
    \mathbf{U}_k + \Delta \mathbf{U}_k.
\]

\subsection{Пошаговое задание смещения}

Для повышения устойчивости нелинейного решения заданное смещение
прикладывается постепенно. На шаге \(m\) задается величина
\[
    \delta_x^{(m)}
    =
    \frac{m}{N_{\mathrm{steps}}}\delta_x,
    \qquad
    m=1,\ldots,N_{\mathrm{steps}}.
\]

В базовых расчетах используется
\[
    N_{\mathrm{steps}}=12.
\]

Таким образом, задача решается как последовательность близких
равновесных состояний, что повышает устойчивость метода Ньютона при
больших деформациях.

\subsection{Вычисление силы реакции}

После нахождения поля перемещений вычисляется сила реакции на верхней
управляющей грани. Вектор поверхностной тяги в лагранжевой постановке
задается выражением
\[
    \mathbf{t} = \mathbf{P}\mathbf{N}.
\]

Компонента реакции по оси \(x\) вычисляется интегрированием по верхней
грани:
\[
    F_x
    =
    \int_{\Gamma_{\mathrm{top}}}
    (\mathbf{P}\mathbf{N})\cdot \mathbf{e}_x\,dS,
\]
где
\[
    \mathbf{e}_x =
    \begin{pmatrix}
        1\\0\\0
    \end{pmatrix}.
\]

В численной реализации используется модуль этой величины:
\[
    \boxed{
    F_x
    =
    \left|
    \int_{\Gamma_{\mathrm{top}}}
    (\mathbf{P}\mathbf{N})\cdot \mathbf{e}_x\,dS
    \right|.
    }
\]

Данная величина интерпретируется как сила, необходимая для поддержания
заданного горизонтального смещения верхней грани реснички.

\subsection{Расчетные серии}

В актуальной версии численной модели выполняются несколько серий расчетов.

Первая серия предназначена для построения облегченной калибровочной
кривой
\[
    F_x = F_x(\delta_x)
\]
при фиксированной сетке \(h=50~\mu\text{м}\), элементах \(P_1\) и
коэффициенте Пуассона \(\nu=0.49\). Рассматриваются значения смещения:
\[
    \delta_x =
    0.25,\;0.30,\;0.35,\;0.40~\text{мм}.
\]

Вторая серия предназначена для грубой проверки сходимости по сетке при
фиксированном смещении
\[
    \delta_x = 0.40~\text{мм}.
\]
Используются размеры сетки
\[
    h=60~\mu\text{м},
    \qquad
    h=50~\mu\text{м}.
\]

Третья серия является диагностикой объемного запирания. Для этого
выполняется расчет при меньшем коэффициенте Пуассона:
\[
    \nu=0.45.
\]
Если сила реакции существенно отличается от результата при
\(\nu=0.49\), это указывает на чувствительность модели к почти
несжимаемости материала и возможное влияние locking-эффекта.

Четвертая серия предназначена для оценки влияния порядка конечного
элемента. Выполняется расчет с квадратичными элементами \(P_2\) при
\[
    h=60~\mu\text{м},
    \qquad
    \nu=0.49,
    \qquad
    \delta_x=0.40~\text{мм}.
\]

\subsection{Итоговая механическая модель}

Итак, в актуальной постановке решается нелинейная трехмерная задача
гиперупругого деформирования двухслойной цилиндрической реснички на
подложке. Неизвестной является вектор перемещений \(\mathbf{u}\).
Материал описывается сжимаемой моделью Neo-Hookean, параметры которой
зависят от слоя. Нижняя поверхность подложки жестко закреплена, а на
всей верхней грани реснички задается горизонтальное смещение
\[
    u_x=\delta_x.
\]

Основным результатом механического расчета является поле перемещений
\[
    \mathbf{u}_h(\mathbf{X})
\]
и сила реакции
\[
    F_x(\delta_x),
\]
которая характеризует механическую жесткость чувствительного элемента.
Полученное деформированное состояние далее используется как основа для
расчета магнитного поля, создаваемого магнитными включениями верхнего
композитного слоя.


% ============================================================
% ГЛАВА 4
% ============================================================

\section{Калибровка, валидация и выбор финальной конфигурации модели}

После построения трёхмерной гиперупругой конечно-элементной модели двухслойной магнитной реснички была проведена серия калибровочных и валидационных расчётов. Основной целью данного этапа являлся подбор такой конфигурации модели, которая обеспечивала бы физически корректное описание деформации и воспроизводила требуемую реакцию конструкции порядка
\[
F_x \approx 60\,\mu\text{Н}.
\]

\subsection{Постановка задачи калибровки}

В механической модели использовалась displacement-control постановка. На верхнюю грань реснички накладывалось кинематическое граничное условие:
\[
u_x = \delta_x,
\]
где $\delta_x$ --- задаваемое горизонтальное смещение верхней поверхности.

При этом вычисляемой величиной являлась реакция конструкции:
\[
F_x
=
\int_{\Gamma_{\text{top}}}
\mathbf{P}\mathbf{N}\cdot \mathbf{e}_x
\, dS,
\]
где
$\mathbf{P}$ --- первый тензор напряжений Пиолы--Кирхгофа,
$\mathbf{N}$ --- внешняя нормаль к верхней поверхности,
$\mathbf{e}_x$ --- единичный вектор оси $x$.

Таким образом, задача калибровки сводилась к определению такого значения $\delta_x$, при котором вычисленная реакция конструкции совпадала бы с целевым значением:
\[
F_x \approx 60\,\mu\text{Н}.
\]

\subsection{Первичная калибровка модели}

На первом этапе использовалась равномерная тетраэдральная сетка с характерным размером элемента:
\[
h = 60\,\mu\text{м}.
\]

Для различных значений управляющего смещения были выполнены нелинейные гиперупругие расчёты.

\begin{table}[H]
\centering
\caption{Первичная калибровка модели при равномерной сетке}
\begin{tabular}{|c|c|}
\hline
$\delta_x$, мм & $F_x$, $\mu$Н \\
\hline
0.20 & 6.005 \\
0.25 & 11.176 \\
0.30 & 18.038 \\
0.32 & 21.216 \\
0.35 & 26.469 \\
0.40 & 36.357 \\
\hline
\end{tabular}
\end{table}

Полученные результаты показали, что диапазон
\[
\delta_x = 0.20\text{--}0.40\ \text{мм}
\]
не позволяет достичь требуемой реакции порядка
\[
60\,\mu\text{Н}.
\]

Однако главным результатом данного этапа стало выявление существенной чувствительности модели к размеру конечных элементов.

\subsection{Проверка сходимости по сетке}

Для проверки устойчивости численного решения была проведена серия расчётов при различных размерах сетки. Для одного и того же значения смещения
\[
\delta_x = 0.32\ \text{мм}
\]
были получены следующие значения реакции конструкции:

\begin{table}[H]
\centering
\caption{Проверка сходимости по сетке для равномерной дискретизации}
\begin{tabular}{|c|c|}
\hline
$h$, $\mu$м & $F_x$, $\mu$Н \\
\hline
100 & 11.854 \\
80 & 7.821 \\
70 & 8.538 \\
60 & 21.216 \\
\hline
\end{tabular}
\end{table}

Данные результаты показали отсутствие устойчивой сеточной сходимости. Вычисленная реакция конструкции существенно изменялась при изменении размера конечных элементов.

Основной причиной являлось недостаточное разрешение поперечного сечения реснички. При диаметре
\[
D = 120\,\mu\text{м}
\]
и равномерной сетке
\[
h = 60\,\mu\text{м}
\]
по диаметру реснички располагалось лишь:
\[
\frac{D}{h}=2
\]
конечных элемента.

Для задачи изгиба тонкой цилиндрической структуры такое разрешение оказалось недостаточным.

\subsection{Переход к раздельной пространственной дискретизации}

Для улучшения качества расчёта была введена неравномерная пространственная дискретизация модели.

Использовались два независимых параметра сетки:
\[
h_{\text{cil}}
\]
--- размер элемента в области реснички,
\[
h_{\text{sub}}
\]
--- размер элемента в подложке.

Подобный подход позволил существенно уточнить дискретизацию в области изгиба реснички без значительного увеличения общего числа конечных элементов.

На промежуточном этапе использовались параметры:
\[
h_{\text{cil}} = 30\,\mu\text{м},
\qquad
h_{\text{sub}} = 100\,\mu\text{м}.
\]

При данных параметрах были выполнены повторные калибровочные расчёты.

\begin{table}[H]
\centering
\caption{Калибровка модели при раздельной сетке}
\begin{tabular}{|c|c|}
\hline
$\delta_x$, мм & $F_x$, $\mu$Н \\
\hline
0.20 & 6.935 \\
0.25 & 12.342 \\
0.30 & 19.447 \\
0.32 & 22.746 \\
0.35 & 28.157 \\
0.40 & 38.344 \\
\hline
\end{tabular}
\end{table}

Результаты показали более физически согласованное поведение модели. При уточнении сетки вычисленная реакция начала изменяться монотонно, что свидетельствовало о приближении к режиму сеточной сходимости.

\subsection{Финальная сеточная валидация}

Для дальнейшего повышения точности расчёта размер элемента в области реснички был уменьшен до:
\[
h_{\text{cil}} = 20\,\mu\text{м}.
\]

При этом размер элемента в подложке оставался:
\[
h_{\text{sub}} = 100\,\mu\text{м}.
\]

Таким образом, по диаметру реснички обеспечивалось:
\[
\frac{D}{h_{\text{cil}}}=6
\]
элементов, что существенно улучшило описание изгибной деформации.

Для проверки влияния смещения верхней грани были выполнены расчёты:

\[
\delta_x = 0.50\ \text{мм}
\]
и
\[
\delta_x = 0.62\ \text{мм}.
\]

Полученные реакции составили:

\[
F_x = 48.793\,\mu\text{Н},
\]
и
\[
F_x = 75.953\,\mu\text{Н}.
\]

На основании линейной интерполяции между данными значениями был выбран следующий кандидат:
\[
\delta_x \approx 0.55\ \text{мм}.
\]

\subsection{Финальный расчёт}

Финальный нелинейный гиперупругий расчёт был выполнен при следующих параметрах:

\[
\delta_x = 0.55\ \text{мм},
\]
\[
h_{\text{cil}} = 20\,\mu\text{м},
\]
\[
h_{\text{sub}} = 100\,\mu\text{м},
\]
\[
\nu = 0.49.
\]

В результате была получена реакция конструкции:
\[
F_x = 59.635\,\mu\text{Н},
\]
что отличается от целевого значения
\[
60\,\mu\text{Н}
\]
менее чем на
\[
1\%.
\]

Дополнительно контролировалась физическая корректность деформации посредством анализа якобиана преобразования:
\[
J = \det \mathbf{F}.
\]

В финальном расчёте было получено:
\[
J \in [0.9993,\ 1.0005].
\]

Полученные значения подтверждают отсутствие вырождения конечных элементов и физическую допустимость деформации.

Максимальное эквивалентное напряжение по Мизесу составило:
\[
\sigma_{\text{vm}}^{\max}
=
2.16\times10^5\ \text{Па}.
\]

Таким образом, в качестве итоговой конфигурации механического расчёта была выбрана модель со следующими параметрами:
\[
\boxed{
\delta_x = 0.55\ \text{мм},
\qquad
h_{\text{cil}} = 20\,\mu\text{м},
\qquad
h_{\text{sub}} = 100\,\mu\text{м}.
}
\]

Данная конфигурация обеспечивает реакцию конструкции, близкую к требуемому значению, и демонстрирует устойчивое и физически корректное поведение гиперупругой конечно-элементной модели.


% ============================================================
% СПИСОК ИСТОЧНИКОВ
% ============================================================
\newpage
\begin{thebibliography}{99}
\addcontentsline{toc}{section}{Список использованных источников}

\bibitem{ribeiro2017}
Ribeiro P., Cardoso S., Bernardino A., Jamone L.
\newblock Bio-inspired ciliary force sensor for robotic platforms.
\newblock \textit{IEEE Robotics and Automation Letters}, 2017, Vol. 2, No. 2, pp. 971--976.

\bibitem{man2023}
Man J., Zhang J., Chen G., Xue N., Chen J.
\newblock A tactile and airflow motion sensor based on flexible double-layer magnetic cilia.
\newblock \textit{Microsystems \& Nanoengineering}, 2023, Vol. 9, Article 12.
\newblock DOI: \url{https://doi.org/10.1038/s41378-022-00478-9}

\bibitem{Zienkiewicz2005}
Zienkiewicz O. C., Taylor R. L., Zhu J. Z.
The Finite Element Method: Its Basis and Fundamentals.
Elsevier, 2005.

\bibitem{daw2019}
Daw A., Pender J.
On the distributions of infinite server queues with batch arrivals //
Journal of Applied Probability. 2019. Vol. 56, no. 1. P. 1--34.

\bibitem{courant1943}
Courant R.
\newblock Variational methods for the solution of problems of equilibrium and vibrations.
\newblock {\em Bulletin of the American Mathematical Society}, 49(1):1--23, 1943.
\newblock Available at: \url{https://projecteuclid.org/journals/bulletin-of-the-american-mathematical-society/volume-49/issue-1/Variational-methods-for-the-solution-of-problems-of-equilibrium-and/bams/1183504922.pdf}

\bibitem{turner1956}
Turner M. J., Clough R. W., Martin H. C., Topp L. J.
\newblock Stiffness and deflection analysis of complex structures.
\newblock {\em Journal of the Aeronautical Sciences}, 23(9):805--823, 1956.
\newblock Available at: \url{https://www.ce.memphis.edu/7117/notes/presentations/papers/Turner%20et%20al%20(1956)%20Stiffness%20and%20deflection%20analysis%20of%20comlex%20strucutres.pdf}

\bibitem{brenner2008}
Brenner S. C., Scott R.
\newblock The Mathematical Theory of Finite Element Methods.
\newblock Springer, 2008.


\bibitem{zhao2019}
Zhao R., Kim Y., Chester S. A., Sharma P., Zhao X.
\newblock Mechanics of hard-magnetic soft materials.
\newblock \textit{Journal of the Mechanics and Physics of Solids}, 2019, Vol. 124, pp. 244--263.

\bibitem{kim2022}
Kim Y., Yuk H., Zhao R., Chester S. A., Zhao X.
\newblock Magnetic soft materials and robots.
\newblock \textit{Chemical Reviews}, 2022, Vol. 122, No. 5, pp. 5317--5364.

\bibitem{Bossavit1998}
Bossavit A. Computational Electromagnetism: Variational Formulations,
Complementarity, Edge Elements. San Diego: Academic Press, 1998.

\bibitem{Jin2014}
Jin J.-M. The Finite Element Method in Electromagnetics. 3rd ed.
Wiley-IEEE Press, 2014.

\bibitem{Kuczmann2008}
Kuczmann M., Ivanyi A. The Finite Element Method in Magnetics.
Budapest: Akademiai Kiado, 2008.

\bibitem{Bermudez2008}
Bermudez A., Rodriguez R., Salgado P.
A finite element method for the magnetostatic problem in terms of scalar potentials.
SIAM Journal on Numerical Analysis, 2008.

\bibitem{BiroPreis1989}
Biro O., Preis K.
On the use of the magnetic vector potential in the finite element analysis of
three-dimensional eddy currents.
IEEE Transactions on Magnetics, 1989, Vol. 25, pp. 3145--3159.

\bibitem{Wang2016}
Wang H., de Boer G., Kow J., Alazmani A., Ghajari M., Hewson R., Culmer P.
Design methodology for magnetic field-based soft tri-axis tactile sensors.
Sensors, 2016, Vol. 16, Article 1356.

\bibitem{DeBoer2017}
de Boer G., Raske N., Wang H., Kow J., Alazmani A., Ghajari M.,
Culmer P., Hewson R.
Design optimisation of a magnetic field based soft tactile sensor.
Sensors, 2017, Vol. 17, Article 2539.

\bibitem{Rosle2019}
Rosle M. H., Wang Z., Hirai S.
Geometry optimisation of a Hall-effect-based soft fingertip for estimating
orientation of thin rectangular objects.
Sensors, 2019, Vol. 19, Article 4056.



\end{thebibliography}

\newpage

% ============================================================
% ПРИЛОЖЕНИЯ
% ============================================================

\appendix

\section*{Приложение А}
\addcontentsline{toc}{section}{Приложение А}

Дополнительные материалы, таблицы, графики или программный код.

\end{document}